% ---------- 自序 ----------
\chapter*{自序：一个没法用三句话回答的问题}
\addcontentsline{toc}{chapter}{自序：一个没法用三句话回答的问题}
\markboth{自序}{自序}

去年冬天，一个做平面设计的朋友发微信问我一件事。

她说她妈把攒了二十年的钱存了三年定期，到期一看，利息比上一回少了一大截。她妈很生气，认定是银行坑人。朋友问我：到底是不是银行坑人。

我打了三行字，删掉。又打了五行，又删掉。

因为要回答这个问题，得先说清楚存款利率是怎么定出来的；要说清楚这个，得先说银行靠什么赚钱；要说这个，绕不开净息差；要说净息差为什么被压到 1.41\%，得说贷款需求为什么起不来；而这又牵到房地产的长周期、牵到央行的政策利率、牵到整个全球的利率环境。

任何一环跳过去，答案就变成了「你就信我」。

那天晚上我意识到一件事：这不是一个能用三句话回答的问题。而普通人关于钱的疑问，几乎全都是这样。它们看上去是一个个孤立的困惑，实际上挂在同一台机器上，而那台机器从来没有人完整地拆给他们看过。

\bookbreak

市面上关于钱的书大致分两类。

一类教你怎么赚钱。这类书必须提供确定性，否则卖不动：这个方法有效，这个位置该买，这个比例是对的。

另一类是正经的教科书。它们是对的，写得也严谨，但它们默认你已经有理由读下去。第一章通常从「货币的三大职能」开始，而一个想弄明白自己家存款为什么缩水的人，读到第三页就合上了。

中间那一块是空的。

一本从零讲起、把机器拆开摆给你看、承认自己不知道很多事、并且明确不打算让你发财的书。

这本书想填的就是中间那块。

\bookbreak

写的时候给自己定了五条约束，现在写在前面，方便你随时拿它们来验收。

\begin{flowlist}
\flowitem{01}{只讲一条链，不做名词表}{名词随时能查，因果链没人讲。全书从头到尾只讲一条链：钱是记账，记账要有人信，银行放贷时创造存款，所以钱约等于债，债要付息，付息需要增长，增长需要投资，投资需要给风险定价，定价错误加杠杆等于危机。二十九篇都挂在这条链的某一环上。}
\flowitem{02}{每个术语先讲人话，再给名字}{先说清它解决什么问题、在生活里长什么样，最后才告诉你它叫什么。倒过来的顺序是教科书的顺序，不是理解的顺序。}
\flowitem{03}{所有数字都标时间}{金融世界的数字变得很快，结构变得很慢。数字请当水位刻度读，结构请当河道读。}
\flowitem{04}{不预测}{书里没有一句关于明年涨跌的判断，没有一只被推荐的标的。这不是谨慎，是这本书的立场，理由写在后记二里。}
\flowitem{05}{不知道就说不知道}{全书至少有五处明确写着「目前没人知道答案」。那五处不是偷懒，恰恰是最该被记住的地方。因为骗你的人从不说这句话。}
\end{flowlist}

\bookbreak

还有一件事得说在前面。

这本书写于 2026 年 8 月。这个时点有点特别：负利率时代刚刚结束，日本央行的政策利率回到 1.0\%，是 1995 年以来的最高；黄金在各国央行的储备里第一次超过了美国国债；一轮由人工智能驱动的资本开支正在改写全球的资金流向；而中国这边，房地产还在长周期的调整里，商业银行的净息差压在历史低位。

三年之后回头看，这本书里的每一个数字都会不对。

但那台机器还是那台机器。银行创造货币的方式不会变，利率传导的路径不会变，泡沫的五幕剧本从 1637 年演到今天也没换过。这本书真正想交给你的，是那些不会变的部分。

\bookbreak

最后。

我那位朋友的问题，后来我用了一顿饭的时间讲完。她听到一半的时候放下筷子，说了一句：

「原来不是银行坑人，是整个环境变了。」

这本书要做的就是这件事。只不过讲给更多人听，而且讲得更慢一点。

\vspace{1em}
\noindent\hfill{\kaishu\color{inkgray} 2026 年 8 月}

\cleardoublepage
